\chapter{Introduction}

Travel introduces unique security risks to your digital
assets and sensitive information. Proper preparation
before, vigilance during, and careful review after travel
creates a comprehensive defense strategy that balances
safety with practical usability.

When we travel, our normal security routines are
disrupted, and we face elevated risks from physical theft,
digital surveillance, border inspections, and social
engineering.

Web3 professionals face additional challenges when
traveling with crypto assets, access to treasury funds, or
sensitive repositories.

The resources in this mini-book provide practical
guidance for maintaining operational security while
traveling.

You can access the contents of this book, and a possibly
more updated version, at
\url{https://frameworks.securityalliance.org}.

\section{Disclaimer}

This is not, by any means, an extensive guide on this
subject or expected to be followed to the T. Its intention
is to guide and provide hints as to where to apply
security. This will vary depending on case to case, or, in
other words, the risks you expose yourself to, by
specifically traveling.

The contents are free, public and open to anyone at the
Security Frameworks official website. Feel free to suggest
an update by contributing directly there.

\section{Three-phase approach}

This travel security mini-book is organized into three
critical phases:

\begin{enumerate}
  \item \textbf{Pre-travel preparation:} Risk assessment,
  device hardening, backup creation, and contingency
  planning.
  \item \textbf{On-trip vigilance:} Network security,
  physical device protection, social engineering
  awareness, and maintaining operational security.
  \item \textbf{Post-travel review:} Device inspection,
  account security verification, and lessons learned
  documentation.
\end{enumerate}

Additional considerations are provided for high-risk
travelers who may face targeted threats due to their role
or access to valuable assets.
